\clearpage\thispagestyle{empty}
\pdfbookmark[0]{Quellen zu Instrumente (Jahrbuch für 1822)}{instruments1822-sources}
\begin{center}\textbf{Quellen zu Instrumente (Jahrbuch für 1822)}\end{center}
\noindent Diese KI-Übertragung ins heutige Deutsch folgt \textit{Instrumente} in den \href{https://gdz.sub.uni-goettingen.de/id/PPN236007467}{\textit{Werken}, Band~VI} (Göttingen, 1874), S.~424--428. Im von Johann Elert Bode herausgegebenen und 1819 in Berlin erschienenen \textit{Astronomischen Jahrbuch für das Jahr 1822} nennt die \href{https://api.digitale-sammlungen.de/iiif/image/v2/bsb10538351_00241/full/1600,/0/default.jpg}{Überschrift des Erstdrucks} Gauß als Verfasser der Beobachtungen und der Instrumentenbeschreibung und Dirksen als Urheber der Bahnelemente des Kometen. Im Exemplar der Bayerischen Staatsbibliothek steht der Bericht auf S.~235--241. Gauß schreibt in der ersten Person und spricht den Empfänger direkt an.
\par\smallskip\noindent In den \href{https://api.digitale-sammlungen.de/iiif/image/v2/bsb10538351_00242/full/1600,/0/default.jpg}{Bahnelementen des Erstdrucks} ist die Zeit der Sonnennähe am 27. Juni mit einem hochgestellten h angegeben. Die Übertragung folgt der Werkausgabe mit dem 28. Juni und hochgestelltem n. Sie behält die unvollständige Wendung über die Genauigkeit der Zenitdistanzen sowie beide Aussagen über den Wert von \(a\) bei. Die \href{https://api.digitale-sammlungen.de/iiif/image/v2/bsb10538351_00246/full/1600,/0/default.jpg}{historische Fußnote zu den beiden Gleichungen} steht in beiden Ausgaben. Der Verweis auf die Gewichte führt in der Werkausgabe zum \hyperlink{native-415}{früheren Bericht über den Meridiankreis von Repsold}. KI-Übertragung und KI-Satz.
\clearpage
